\subsection{Deep World Models}
\label{section:fc_ha_schmidhuber}

\citet{HaSchmidhuber2018} trains a visual encoder, a recurrent latent-state model, and a controller on pixel observations. For CarRacing, the controller is trained and evaluated in the actual environment. For VizDoom, the controller is trained in the learned recurrent model and then evaluated in the actual environment. The system has three components with separate objectives. The vision model reconstructs observations, the memory model predicts the next latent state, and the controller maximizes returns. The paper draws directly on the controller-model architecture of \citet{Schmidhuber1990}. Its implementation combines a convolutional variational autoencoder, a mixture-density recurrent network, and an evolution strategy to operate on pixel observations.

A \emph{simulated rollout} is a trajectory generated by a learned model, and a \emph{planning update} uses model-generated experience to update a policy or value function. Neural world models often encode observations and histories in a compressed \emph{latent state}. Training minimizes a \emph{model loss}; a \emph{value-aware loss} penalizes errors according to their effect on value estimates. A \emph{short rollout} limits the number of model-generated transitions before a \emph{bootstrap} supplies an estimated value for the remaining return. A \emph{model error} is a discrepancy between the learned model and the environment, and its importance for control depends on how it affects planning.

\subsubsection{Three components}
\label{section:fc_ha_schmidhuber_components}

The vision component $V$ is a convolutional variational autoencoder.\footnote{A variational autoencoder compresses high-dimensional observations into a low-dimensional latent vector by training an encoder and a decoder jointly.} The encoder maps an observation $x_t \in \mathbb{R}^{64 \times 64 \times 3}$ to a latent vector $z_t \sim \mathcal{N}(\mu(x_t), \sigma(x_t)^2 I)$ with $z_t \in \mathbb{R}^{N_z}$, where $N_z = 32$ for CarRacing and $N_z = 64$ for VizDoom. The decoder is trained jointly with the encoder against the standard ELBO objective,\footnote{Evidence lower bound; the standard variational autoencoder training objective that combines a reconstruction term with a Kullback-Leibler penalty between the encoder's latent and a fixed prior. The Kullback-Leibler term measures the discrepancy between two distributions.} combining a per-pixel reconstruction term with a Kullback-Leibler penalty on the latent. The vision component is trained on observation reconstruction alone, independently of the controller, on a buffer of frames collected by a random policy. After training, the decoder is not used by $M$ or $C$, which operate on the encoder's latent vectors rather than pixels. The lower-dimensional representation reduces the cost of controller optimization and model rollouts.

The memory component $M$ is a mixture-density recurrent neural network.\footnote{A mixture-density recurrent network is a recurrent network whose output parameterizes a Gaussian mixture over the next latent rather than a single mean; this matters when the next state is genuinely multimodal.} Its body is an LSTM with hidden state $h_t$.\footnote{Long short-term memory; a gated recurrent cell that propagates information across time and is the dominant pre-Transformer architecture for sequential data.} Its head outputs the parameters of a Gaussian mixture over the next latent, conditional on the current action, the current latent, and the recurrent state, modeling $P(z_{t+1} \mid a_t, z_t, h_t)$ as a mixture of five diagonal Gaussians. For controller training inside the learned model, a temperature parameter changes the distribution from which the next latent state is sampled. The memory component is trained on next-latent likelihood alone, again independently of the controller, using the same buffer of rollouts under a random policy that trained $V$. The mixture can represent discrete events such as whether a monster fires in VizDoom. A unimodal Gaussian may instead average across the possible outcomes.

The controller $C$ is a single linear layer. It outputs an action $a_t = W_c [z_t, h_t] + b_c$ from the concatenation of the visual latent and the recurrent state. On CarRacing the controller has roughly $870$ parameters, compared to about $5$ million in $V$ and $M$ together. It is trained by a covariance matrix adaptation evolution strategy,\footnote{A gradient-free black-box optimizer that maintains a Gaussian distribution over candidate parameter vectors and updates that distribution from their evaluated returns.} CarRacing candidates are evaluated in the actual environment, while VizDoom candidates are evaluated in the learned latent environment. The controller is the only component whose objective depends on return.

\subsubsection{Empirical Results}
\label{section:fc_ha_schmidhuber_empirics}

On CarRacing-v0 the full agent scores $906 \pm 21$ over $100$ random tracks. The benchmark requires an average above $900$ to count as solved, and this is the first reported solution. On the VizDoom DoomTakeCover task the controller is trained inside the memory network, with the visual encoder and recurrent forecaster fitted on a buffer of random rollouts and the controller never exposed to real frames during policy search. Transferred zero-shot to the real environment, the controller scores $1092 \pm 556$ against a solution threshold of $750$. The variance is large but the mean clears the threshold, and the policy is constructed without any direct contact with the real environment after the data collection phase.

\subsubsection{Training Schedule and Limitations}
\label{section:fc_ha_schmidhuber_reading}

The authors train $V$, $M$, and $C$ separately, reporting that this is more practical than end-to-end training. Each component has a separate objective and optimizer, and training proceeds sequentially rather than jointly. This schedule does not train $V$ for task relevance, and the dynamics objective is independent of controller return.

\subsection{The Recurrent State-Space Model and the Dreamer Line}
\label{section:fc_rssm}

PlaNet and Dreamer replace the separately trained vision and memory components of \citet{HaSchmidhuber2018} with a jointly trained recurrent state-space model. PlaNet uses this model for online planning, while DreamerV1 trains an actor and critic on imagined latent trajectories. DreamerV2 changes the stochastic latent variables and Kullback-Leibler objective, and DreamerV3 adds mechanisms for robustness across tasks with different signal scales.

\subsubsection{The Recurrent State-Space Model}
\label{section:fc_rssm_transition}

Let $o_t$ denote the observation at time $t$, $a_t$ the action, and $r_t$ the reward. The Recurrent State-Space Model factors the latent state into a deterministic recurrent component $h_t$ and a stochastic latent $s_t$, coupled by
\[ h_t = f_\theta(h_{t-1}, s_{t-1}, a_{t-1}), \qquad s_t \sim p_\theta(s_t \mid h_t). \]
The deterministic path $f_\theta$ is a recurrent cell, implemented as a Gated Recurrent Unit in PlaNet. PlaNet and DreamerV1 use a diagonal Gaussian prior, whereas DreamerV2 and DreamerV3 use vectors of categorical latent variables. A posterior network $q_\theta(s_t \mid h_t, o_t)$ incorporates the current observation when inferring the latent state from observed trajectories. A simplified length-$T$ objective for PlaNet and DreamerV1 contains three terms,
\[ \mathcal{L}_\theta = \underbrace{-\sum_t \log p_\theta(o_t \mid h_t, s_t)}_{\text{reconstruction}} \;-\; \underbrace{\sum_t \log p_\theta(r_t \mid h_t, s_t)}_{\text{reward log-likelihood}} \;+\; \underbrace{\sum_t \mathrm{KL}\!\left[ q_\theta(s_t \mid h_t, o_t) \,\Vert\, p_\theta(s_t \mid h_t) \right]}_{\text{prior-posterior consistency}}. \]
DreamerV1 predicts the discount factor on tasks with early termination, and DreamerV2 and DreamerV3 include a discount or continuation predictor in the world model. DreamerV3 separates the Kullback-Leibler term into dynamics and representation losses. PlaNet proposes latent overshooting for training multi-step latent predictions, although its final RSSM agent does not require it. The recurrent state carries information across time, while the stochastic state represents multiple possible futures. PlaNet evaluates candidate action sequences with the learned model, whereas Dreamer trains an actor and critic on model-generated trajectories.

\subsubsection{Actor-critic in imagination}
\label{section:fc_rssm_imagination}

Dreamer trains an actor $\pi_\phi(a_t \mid h_t, s_t)$ and a critic $V_\psi(h_t, s_t)$ on trajectories generated by the learned latent dynamics. These trajectories begin from posterior model states inferred from replayed observations. A schematic finite-horizon objective is
\[ J(\phi) = \mathbb{E}_{\pi_\phi, p_\theta}\!\left[ \sum_{\tau = t}^{t+H} \gamma^{\tau - t} \hat r_\tau + \gamma^{H+1} V_\psi(h_{t+H+1}, s_{t+H+1}) \right], \]
Here $H$ is the imagination horizon, $\hat r_\tau$ is a predicted reward, and $V_\psi$ bootstraps rewards beyond that horizon. DreamerV1 backpropagates actor gradients through reparameterized actions and latent states in continuous control and uses straight-through gradients for categorical actions. DreamerV2 uses REINFORCE for Atari and dynamics backpropagation for continuous control, while DreamerV3 uses REINFORCE for both action types. Each Dreamer generation trains the critic on $\lambda$-return targets from imagined trajectories.

Like Dyna-Q, Dreamer uses a learned model to improve behavior without additional environment interaction. The update differs because Dyna-Q applies tabular Q-learning to sampled one-step transitions, whereas Dreamer trains an actor and critic on multi-step latent trajectories. These nonlinear actor-critic updates do not have the tabular convergence guarantee available for Dyna-Q. Theorem~2 of \citet{Asadi2018Lipschitz}, discussed in \S\ref{section:fc_synthesis_fails}, gives a different result for a Markov reward process with a Lipschitz reward function. If the one-step model is $\Delta$-accurate in Wasserstein distance and $\bar K = \min\{K_F,K_T\}$ satisfies $\gamma \bar K < 1$, it bounds the resulting value error.\footnote{PlaNet \citep{HafnerPlaNet2019} uses an RSSM with online cross-entropy planning. DreamerV1 \citep{HafnerDreamer2020} replaces online planning with actor-critic training on imagined RSSM trajectories. DreamerV2 \citep{HafnerDreamerV2_2021} replaces Gaussian latent variables with vectors of categoricals and adds Kullback-Leibler balancing. Under its 200M-step Atari protocol, it exceeds IQN and Rainbow under all four reported aggregate metrics. DreamerV3 \citep{HafnerDreamerV3_2025} uses one hyperparameter configuration across more than one hundred and fifty tasks, with symlog transformations, percentile return normalization, and free bits among its robustness mechanisms. It also collects Minecraft diamonds from scratch without human data or curricula.}

PlaNet and Dreamer train their world models with predictive losses for observations, rewards, and latent transitions. The Dreamer versions also predict episode discount or continuation where applicable. These losses measure predictive fit rather than the error in value estimates later computed from the latent state. In economic terms, the RSSM is a high-dimensional perceived law of motion trained with predictive losses rather than hand-chosen moment conditions. Value-aware methods instead penalize model errors according to their effect on value functions \citep{FarahmandVAML2017,GrimmVALEQ2020}. MuZero follows a related empirical design by training latent dynamics with reward, policy, and value targets rather than observation reconstruction.

\subsection{Short Model Rollouts and Dynamics Ensembles}
\label{section:fc_mbpo}

\citet{janner2019model} applies short model rollouts to continuous-control policy optimization. Its implementation uses a probabilistic dynamics ensemble with soft actor-critic. Model rollouts begin from real states sampled from the replay buffer. Starting each rollout from a replay-buffer state allows the model rollout length to differ from the task horizon. PETS uses an ensemble for model-predictive control rather than policy learning \citep{Chua2018PETS}, while ME-TRPO uses ensemble-generated trajectories to train a TRPO policy \citep{Kurutach2018}. MBPO's branching scheme resembles Dyna, which samples previously observed states for one-step simulated updates \citep{sutton1990}. \citet{janner2019model} interprets these as length-one branched rollouts.

\subsubsection{Branched short rollouts}
\label{section:fc_mbpo_branching}

A $k$-step branched rollout begins at a replay-buffer state and follows the current policy in the learned dynamics for $k$ steps. During training, MBPO samples a state from the environment replay buffer, generates a $k$-step rollout under the current policy and learned model, stores the synthetic transitions in a model buffer, and updates soft actor-critic on that buffer.

Theorem~4.2 of \citet{janner2019model} bounds the gap between the true return and the branched-rollout return. With $\epsilon_m$ the model error measured in total variation, $\epsilon_\pi$ the policy distribution shift relative to the data-collection policy, $k$ the branching length, and $r_{\max}$ the reward bound,
\[ \eta[\pi] \ge \eta^{\text{branch}}[\pi] - 2 r_{\max} \Big[ \frac{\gamma^{k+1} \epsilon_\pi}{(1 - \gamma)^2} + \frac{\gamma^k + 2}{1 - \gamma} \epsilon_\pi + \frac{k}{1 - \gamma}(\epsilon_m + 2 \epsilon_\pi) \Big]. \]
The $k$-dependent parts of the first two terms decrease with $k$, while the final term grows linearly in $k$ and depends on both model error and policy shift. Using model error under the updated policy, Theorem~4.3 gives a related penalty that is minimized at some $k^\star>0$ when that error is sufficiently small. In its Hopper ablation, the paper reports that fixing $k=1$ retains much of MBPO's benefit. The reported task-specific settings use either a fixed one-step rollout or a schedule that increases from one step to 15 or 25 steps. In the same Hopper ablation, $200$-step rollouts remain usable but perform worse, while $500$-step rollouts are too inaccurate for effective learning.

MBPO uses an ensemble of seven probabilistic neural networks. Each network predicts a diagonal Gaussian over the next state and reward, and the simulator selects one ensemble member uniformly at each synthetic transition. The Gaussian output represents aleatoric uncertainty, while the bootstrapped ensemble represents epistemic uncertainty.\footnote{\citet{Kurutach2018} trains TRPO on simulated trajectories generated by selecting an ensemble member independently at each transition. It evaluates policy improvement across the ensemble every five updates, uses $70\%$ as the validation threshold, and stops policy training after repeated failures of this criterion. \citet{Chua2018PETS} uses a five-member probabilistic ensemble and particle-based trajectory sampling inside model-predictive control. It optimizes an action sequence at each real step, executes the first action, and replans.}\footnote{On the standard $1{,}000$-step MuJoCo tasks, \citet{janner2019model} reports that MBPO's Ant return after $300{,}000$ real steps matches SAC after $3{,}000{,}000$ steps. \citet{Kurutach2018} reports approximately $100$ times fewer real samples than its model-free baselines on six continuous-control tasks, using horizons of $100$ or $200$ steps.}
